\subsection{IV, seconde partie : domaines, points et supports}
\label{note:iv2-domains}
En 3.1.3, l'idéal considéré au point $x$ est $\mathfrak j_x$, avec la
notation de 3.1.2. En 4.8.8, le changement de scalaires de $v'$ vers
$K$ s'écrit $v'\otimes1_K$. En 4.9.5, l'identification avec le corps
intermédiaire est $K=K_1$ ; c'est alors l'extension $K'/K$ qui est
algébrique. En 5.2.5, l'anneau dont le spectre a les deux points annoncés
est $B$.

L'ensemble des zéros de $g$ en 5.3.1 est $X\setminus X_g$ : la notation
$X_g$ désigne l'ouvert où $g$ est inversible. En 5.6.10, un idéal maximal
de l'algèbre finie $B$ est au-dessus de l'idéal maximal de $A$ ; l'extension
algébrique qui intervient dans la formule des dimensions est donc celle
des \emph{corps résiduels}. En 5.6.11, les idéaux de la chaîne sont contenus
dans l'idéal $\mathfrak K$. En 5.7.4, l'inégalité de coprofondeur se lit
au point $z$ considéré, donc sur $\mathscr F_z$.

Les ouverts utilisés pour prolonger les sections à travers $Z$ contiennent
le complément de $Z$ : en 5.9.8 ils contiennent $V\cap(X-Z)$, et en 5.10.12
ils contiennent $X-Z$. La propriété du germe en 5.9.11 est d'être
$Z$-clos, conformément à 5.9.9. Le recouvrement de 5.10.8 comprend aussi
$U_0=X_0\cap U$. Enfin, pour appliquer 5.11.6(c), le faisceau est
$\mathcal G=(A/\mathfrak q)^\sim|U$ et le fermé est $T=V(\mathfrak J)$,
qui sont les objets désignés dans cette condition.

\subsection{IV, seconde partie : hypothèses algébriques}
\label{note:iv2-flatness}
Dans 6.8.5, les morphismes structuraux d'un produit sur un corps sont
plats : les espaces vectoriels sont plats et la platitude est stable par
changement de base. La fidélité n'est pas nécessaire à cet argument et
ne serait pas assurée pour un facteur vide. En 6.10.5, le nombre $r(n)$
est le \emph{rang} du module libre ; sa longueur dépend aussi de celle
de l'anneau de base. En 6.14.1, l'inclusion utilisée est
$\mathfrak a\cap\mathfrak b\subset\mathfrak a\cap\mathfrak c$.

Dans 6.15.10, les propriétés d'intégralité, de birationalité et de
radicialité concernent le morphisme $l$ et son changement de base $l'$.
En 7.1.8, on choisit d'abord le premier $\mathfrak p'$ au-dessus de
$\mathfrak p$ ; sa contraction à $A'$ est ensuite l'idéal maximal
$\mathfrak mA'$. Le choix au-dessus de ce dernier idéal seul ne fixe pas
la contraction intermédiaire.

Les anneaux de 7.3.13 et 7.6.4 sont semi-locaux : la complétion peut
concerner plusieurs idéaux maximaux. L'hypothèse employée en 7.7.2 porte
sur tout quotient intègre de $A_{\mathfrak p}$ ; c'est à ces anneaux
intègres que s'applique la propriété d'être japonais.

Pour le faisceau sans torsion de 7.8.3, la non-nullité assure que le
germe générique n'est pas nul et que l'ensemble des points associés
contient le point générique. Le faisceau nul a un ensemble de points
associés vide. Enfin, en 7.9.1, un anneau local noethérien normal de
dimension au plus un est soit un corps, soit un anneau de valuation
discrète. Ces deux cas donnent la régularité ; le point générique
appartient au premier cas.

\subsection{IV, 6.15.2 : le localisé de la normalisation}
\label{note:iv2-normalization}
On conserve les notations de l'exemple et son hypothèse que $2$ est
inversible dans $K$. Les relations $t^2=u^2-w^2$ et $v=tu/(u+w)$
donnent l'extension quadratique $K(u,w)(\sqrt{u^2-w^2})$.
Posons $a=t-v$ et $C_2=K[t,v,w]$. Puisque $a=tw/(u+w)$,
on a $u=vw/a$ dans le corps des fractions. En substituant dans
$t^2=u^2-w^2$, on obtient
\[
 t a^2=w^2(2v-t),\qquad t(a^2-w^2)=-2aw^2.
\]
Dans $(C_2)_a$, cette dernière identité implique
\[
 \left(1-\frac{t}{2a}\right)(a^2-w^2)=a^2.
\]
Ainsi $a^2-w^2$ est une unité, et les expressions
$t=-2aw^2/(a^2-w^2)$, $v=t-a$ donnent
\[
 (C_2)_a=K[a,w,a^{-1},(a^2-w^2)^{-1}].
\]
Les identités précédentes donnent aussi $K(a,w)=K(u,v,w)$.
Ce corps est de degré de transcendance deux sur $K$ ; $a,w$ sont donc
algébriquement indépendants. Le membre de droite est un localisé d'un
anneau de polynômes, donc est intégralement clos. C'est précisément le
localisé utilisé dans l'intersection de 6.15.2. La preuve n'exige pas
que l'anneau $C_2$ lui-même soit intégralement clos.

\subsection{IV, seconde partie : repères des résultats invoqués}
\label{note:iv2-refs}
Les définitions de 4.6.21 sont celles de 4.6.17 et 4.6.19 ; le critère
local de réduction géométrique invoqué en 4.6.22 est 4.6.13.
Les formules utilisées en 5.5.3, 5.5.7 et 5.9.4 sont respectivement
5.5.1.1, 5.5.7.1 et 5.9.4.1.
La propriété $(\mathrm P'_{\mathrm I})$ en 7.3.11 renvoie à 7.3.5(iii).
Pour 7.4.4(i), le renvoi est 7.4.3(i) ; pour l'équivalence des conditions
$a)$ et $b)$ en 7.4.5, il est 7.4.4(i).
Les propriétés d'excellence et de localisation utilisées en 7.8.4 et dans 7.9.7 renvoient respectivement à 7.8.3(iv) et 7.8.3(ii).
La référence bibliographique [36] est le travail de S.~Abhyankar,
\emph{Local uniformization on algebraic surfaces over ground fields of
characteristic $p\ne0$}, publié dans les \emph{Annals of Mathematics}.
\par
